% STATUS: written 2026-11-06 (English gate clean; sources read in full as recorded below).
% Sources read in full: Weibel 1994 (section 9.2), Hochschild 1945 (recalled, see ch 1.9).
\section{The first Hochschild cohomology group}\label{sec:cohomology:first-cohomology}

Chapter~\ref{ch:origins} defined the Hochschild complex for an associative
algebra $A$ over a field with coefficients in an
$A$-bimodule $M$~\cite[p.~60]{Hochschild1945}. The definition follows
Hochschild~\cite[p.~60]{Hochschild1945}.

The complex has a one-cocycle and a one-coboundary
\cite[p.~60]{Hochschild1945}. A one-cocycle is a derivation of $A$ into $M$,
and a one-coboundary is an inner
derivation~\cite[p.~60]{Hochschild1945}. So the first group measures how far
the derivations are from being inner.

Weibel gives the modern statement of that measurement
\cite[p.~301]{Weibel1994}. It gives the first group as the derivations modulo
the inner ones, and this book uses that form.

\begin{lemma}[{\cite[Lemma 9.2.1, p.~301]{Weibel1994}}]\label{lem:cohomology:h1-der}
$H^1(R, M) = \Der_k(R, M)/\operatorname{PDer}_k(R, M)$, the derivations of $R$ into
$M$ modulo the principal derivations, where $\operatorname{PDer}$ denotes the
principal, that is inner, derivations.
\end{lemma}

Weibel notes a similarity here~\cite[p.~301]{Weibel1994}. The low-dimensional
Hochschild groups have the same shape as the cohomology of groups, and that is
the situation of Theorem~\ref{thm:cohomology:h1-group} below.

So Lemma~\ref{lem:cohomology:h1-der} gives one half of a story. The other half
is for groups, and Section~\ref{sec:cohomology:group-cohomology} completes it.

The first group also controls the higher
ones~\cite[Theorem~3.1, p.~60]{Hochschild1945}.
Theorem~\ref{thm:origins:hochschild-separable} of Chapter~\ref{ch:origins} gives
the reduction of Hochschild. For every bimodule $M$ it produces a bimodule
$(A:M)$ that raises the degree by one:
\[
  H^m(A,M) \cong H^{m-1}(A,(A:M)) \quad \text{for every } m \ge 2
\]
\cite[Theorem~3.1, p.~60]{Hochschild1945}. Two consequences follow from that
reduction~\cite[Theorem~3.1, p.~60]{Hochschild1945}.

The first consequence is that one group decides all the others, so the groups
of dimension one determine the rest of the
theory~\cite[Theorem~3.1, p.~60]{Hochschild1945}.

The second consequence concerns the algebras where the theory
vanishes~\cite[p.~58 and \S~4]{Hochschild1945}. An algebra is separable exactly
when every derivation of it into every bimodule is
inner~\cite[p.~58 and \S~4]{Hochschild1945}. That is the same as the vanishing
of all the cohomology groups~\cite[p.~58 and \S~4]{Hochschild1945}, and it is
the statement of
Theorem~\ref{thm:origins:hochschild-separable} in Chapter~\ref{ch:origins}.

The identification of Lemma~\ref{lem:cohomology:h1-der} has a recent
generalization~\cite[Definition~2.1, p.~1]{BennettHeilLaubacher2023}. Bennett,
Heil and Laubacher work over a commutative triple, in which the second algebra
is commutative and its map into the first one lands in the
centre~\cite[Definition~2.1, p.~1]{BennettHeilLaubacher2023}. They then define
a \emph{secondary derivation} on that
triple~\cite[Definition~3.1, p.~5]{BennettHeilLaubacher2023}.

\begin{definition}[{\cite[Definition~3.1, p.~5]{BennettHeilLaubacher2023}}]\label{def:cohomology:secondary}
A \emph{secondary derivation} of the commutative triple
$\mathcal{T} = (A, B, \varepsilon)$ with values in $M$ is a $k$-linear map
$D\colon B\otimes A \to M$ with the symmetric bimodule structure
$(\alpha\otimes a)\cdot D(\beta\otimes b) = a\varepsilon(\alpha)D(\beta\otimes b)$
such that
\begin{enumerate}
\item $D(\lambda(\alpha\otimes a) + \mu(\beta\otimes b)) = \lambda D(\alpha\otimes a) + \mu D(\beta\otimes b)$,
\item $D((\alpha\otimes a)(\beta\otimes b)) = a\varepsilon(\alpha)D(\beta\otimes b) + D(\alpha\otimes a)b\varepsilon(\beta)$,
\item $D(\alpha\otimes 1) + D(\alpha\otimes 1) = D(1\otimes\varepsilon(\alpha))$
\end{enumerate}
for all $a, b\in A$, $\alpha, \beta\in B$ and $\lambda, \mu\in k$. The set of
such secondary derivations is denoted $\Der_k(\mathcal{T}, M)$.
\end{definition}

When $B = k$ this reduces to the ordinary derivations of
$A$~\cite[Remark~3.3, p.~5]{BennettHeilLaubacher2023}. The recent theorem then
reads as the analogue of Lemma~\ref{lem:cohomology:h1-der}.

\begin{theorem}[{\cite[Theorem~3.6, pp.~5--6]{BennettHeilLaubacher2023}}]\label{thm:cohomology:secondary}
For a commutative triple $\mathcal{T} = (A, B, \varepsilon)$ one has
$HH^1(A, B, \varepsilon) \cong \Der_k(\mathcal{T}, A^*)$.
\end{theorem}

Theorem~\ref{thm:cohomology:secondary} keeps the shape of
Lemma~\ref{lem:cohomology:h1-der}: the first group is again a set of
derivations, and they now live over $B\otimes A$ rather than over $A$
alone~\cite[Theorem~3.6, pp.~5--6]{BennettHeilLaubacher2023}.

Further work in the cohomology line, held in the collection:
\noindent\emph{The Lie module structure on the Hochschild cohomology groups of monomial algebras with radical square zero}~\cite{LibTheLieModuleStructure2008}.
\noindent\emph{The BV-algebra structure of the Hochschild cohomology of the group ring of cyclic groups of prime order}~\cite{LibTheBvAlgebraStructure2016}.
\noindent\emph{Bounded Hochschild cohomology of Banach algebras with a matrix-like structure}~\cite{LibBoundedHochschildCohomologyOf2006}.
\noindent\emph{Gerstenhaber brackets on Hochschild cohomology of quantum symmetric algebras and their group extensions}~\cite{LibGerstenhaberBracketsOnHochschild2016}.
\noindent\emph{Central extensions of the algebra of formal pseudo-differential symbols via Hochschild cohomology and quadratic symplectic Lie algebras}~\cite{LibCentralExtensionsOfThe2018}.
\noindent\emph{The first module $\sigma$ $\tau$ -cohomology group of triangular Banach algebras of order three}~\cite{LibTheFirstModuleCohomology2018}.
\noindent\emph{The first and second Hochschild cohomology groups of Banach algebras with coefficients in special symmetric bimodules}~\cite{LibTheFirstAndSecond2020}.
\noindent\emph{On the cohomology of the spaces of differential operators acting on skewsymmetric tensor fields or on forms as modules of the Lie algebra of vector fields}~\cite{LibOnTheCohomologyOf2004}.
